\begin{afterword}
\summaryhead

The post starts from an old story. Plato's school defined man as a ``featherless biped''; Diogenes the Cynic is said to have brought in a plucked chicken and called it Plato's man, and the Platonists added ``with broad nails.''

No list of features, the author argues, captures everything humans have in common while fitting all humans and only humans. A rough definition is still useful: ``featherless biped'' picks out some real humans, and from them we learn many more shared traits, such as language, tools, red blood and vulnerability to hemlock. The category becomes a ``similarity cluster,'' and a plucked chicken no longer fits, whatever the definition says. A definition is a clue that leads to the cluster, after which it has done its work. So a dictionary should be read as a book of hints for matching words to similarity clusters, not as a book of Aristotelian class definitions.

\responsehead

The post is short and makes one point: categories are held together by many shared traits, and a definition only has to point the way to them. The point is sound, and the post tells its story with the right caution. The source is Diogenes Laertius, a biographer of the third century AD writing about a Cynic of the fourth century BC, and the post says the chicken ``is said to have'' been produced.

The idea has predecessors that the post does not name. Wittgenstein described categories held together by ``a complicated network of similarities overlapping and criss-crossing'' and called them ``family resemblances.'' Psychologists have studied the same structure as ``prototype effects,'' and the post published the same day, ``Typicality and Asymmetrical Similarity,'' cites that research (Rosch and Lloyd, 1978). So the reader of the book meets it one post later. What the post adds is the account of how a rough definition gets a learner to the cluster, and that account is clear.

The foil is weaker than the thesis. The post says that if ``Aristotelian logic'' were a good model of psychology, the Platonists would have accepted the chicken as a man. But Aristotle's own rules for definitions reject a definition that fits something other than the thing defined. By those rules the chicken refutes the definition, which is what the Platonists concluded. The rule assumes, as the post does, that one can tell a chicken from a man before any definition is settled. So Aristotle's practice supports the post's account more than the parenthesis allows.

In short: a sound point about categories as clusters of shared traits, older than the post, set against a version of Aristotelian logic that Aristotle's own rules for definition do not fit.
\end{afterword}
